%%%%%%%%%%%%%%%%%%%%%%%%%%%%%%%%%%%%%%%%%
% IEEE TAES initial-submission cover letter.
% Adapted from the author-provided coverletter.tex template.
% Original template credits: Chen Yang and Fanchao Chen;
% original template source: LaTeXTemplates.com.
% Scientific claims checked against ../main.tex on 9 October 2026.
%%%%%%%%%%%%%%%%%%%%%%%%%%%%%%%%%%%%%%%%%
\documentclass[10pt]{article}
\ifdefined\pdfminorversion\pdfminorversion=7\fi
\usepackage{charter}
\usepackage[a4paper,top=1in,bottom=1in,left=1in,right=1in]{geometry}
\usepackage{graphicx}
\usepackage{fancyhdr}
\usepackage[hidelinks]{hyperref}
\hypersetup{
  pdftitle={Cover Letter for IEEE Transactions on Aerospace and Electronic Systems},
  pdfauthor={Zheng Tan},
  pdfsubject={Regular Paper submission: required-target orbital inspection}
}

\setlength{\parindent}{0pt}
\setlength{\parskip}{0.7em}

\fancypagestyle{firstpage}{%
  \fancyhf{}
  \renewcommand{\headrulewidth}{0pt}
  \renewcommand{\footrulewidth}{1pt}
}
\fancypagestyle{subsequentpages}{%
  \fancyhf{}
  \renewcommand{\headrulewidth}{1pt}
  \renewcommand{\footrulewidth}{1pt}
}
\AtBeginDocument{\thispagestyle{firstpage}}
\pagestyle{subsequentpages}

\newcommand{\SubmissionDate}{9 October 2026}
\newcommand{\CorrespondingAuthor}{Zheng Tan}
% Submission-contact email supplied by the author on 9 October 2026.
\newcommand{\CorrespondingEmail}{\href{mailto:zheng-uav.tan@connect.polyu.hk}{zheng-uav.tan@connect.polyu.hk}}
\newcommand{\ManuscriptTitle}{Viability-Preserving Rollout Approximate Dynamic Programming for Required-Target Orbital Inspection}

\begin{document}

\includegraphics[width=0.4\textwidth]{Polyu-cover.pdf}

\vspace{-1em}
\rule{\linewidth}{1pt}

\vspace{0.6em}
\hfill\begin{tabular}{r@{}}
  \SubmissionDate\\[0.35em]
  \CorrespondingAuthor\\
  Email: \CorrespondingEmail\\
  Department of Aeronautical and Aviation Engineering\\
  The Hong Kong Polytechnic University\\
  Hung Hom, Kowloon, Hong Kong SAR
\end{tabular}

\vspace{0.5em}
Dear Editor,

We submit our manuscript entitled ``\ManuscriptTitle'' for consideration as a Regular Paper in \textit{IEEE Transactions on Aerospace and Electronic Systems}.

Autonomous orbital inspection must observe every prescribed target while limiting maneuver effort under relative-motion, visibility, and safety constraints. The manuscript develops viability-preserving rollout approximate dynamic programming (ADP) over a directed graph of audited orbital arcs. It couples Hill--Clohessy--Wiltshire transfers with camera-valid observations and task-completing base-policy continuations.

The methodological contribution is to preserve a feasible continuation to all required observations during finite-depth rollout. Safe connecting actions remain available even when they provide no immediate observation gain. Under deterministic finite-graph assumptions and exhaustive prefix evaluation, we prove that repeated rollout preserves base-policy completion and does not increase its graph cost.

Independent route certification distinguishes planner failure from unavailable viewpoints. On the resulting 59-case all-solvable development benchmark, depth-three ADP completes 59/59 missions, compared with 57/59 for one-step ADP and 33/59 for task-aware greedy planning. Its mean graph-cost gap is 0.96\% across 27 cases with proven optima, with shorter median runtimes than the independent mixed-integer reference. A separate fixed-library confirmation shows an 11.05\% reduction in mean maneuver velocity increment relative to one-step-ADP-seeded local search on 43 jointly completed test missions. ROS~2/Gazebo simulation execution further demonstrates completion of all nine required targets using twelve observations. These findings concern the specified candidate graphs and enabled deterministic audits.

The work addresses aerospace guidance, autonomous mission planning, and on-orbit servicing, aligning with the journal's focus on aerospace systems. Supporting analyses and simulation execution details are provided in the accompanying supplementary material.

Thank you for your consideration.

Sincerely yours,

\vspace{0.5em}
\CorrespondingAuthor\\
Corresponding author\\
The Hong Kong Polytechnic University

\end{document}
